\documentclass[a4paper,oneside]{article}
\usepackage{phd}
\usepackage{hyperref}
\input{defaultstyle}
\newfontfamily\calibri{calibri.ttf}
\begin{document}
\calibri
\lipsum[2]

\bgroup\parindent=0pt
\leavevmode\par
\includegraphics[width=\linewidth]{./images/pierced.png}%
\vskip0sp %
\ifodd\thepage
   %\nobreak
   \hbox to \linewidth{\hfill\scriptsize Credit:\thinspace{\small\itshape Kathleen Gilje}}%
\else
  \nobreak
  \hbox to \linewidth{\scriptsize Credit:\thinspace{\small\itshape Kathleen Gilje}\hfill}%
\fi     
\enlargethispage{30pt}
\nobreak 
\captionof{figure}{Some description of the image.}%
\bigskip
\egroup

\lipsum[1]

\bgroup\parindent=0pt
\leavevmode\par
\includegraphics[width=\linewidth]{./images/pierced.png}%
\nobreak
\vskip0sp %
\ifodd\thepage
   \hbox to \linewidth{\hfill\scriptsize Credit:\thinspace{\small\itshape Kathleen Gilje}}%
\else
  \hbox to \linewidth{\scriptsize Credit:\thinspace{\small\itshape Kathleen Gilje}\hfill}%
\fi      
\captionof{figure}{Some description of the image.}%
\bigskip
\egroup

\lipsum[1]

\parindent0pt

\tikzrule

\bgroup
\leftskip-3cm \rightskip2cm
\def\hook{\node[right=5pt, yshift=-12pt] at (0,-3) {\HUGE\color{purple} This is the  Title}; }
\def\hook{}
\makeatletter\@specialtrue 
\cxset{chapter name = CHAPTER}
\expandafter\ifnum\thechapter=0\stepcounter{chapter}\else\fi
\begin{tikzpicture}
%\draw [help lines] (0,0) grid (18,-13);
\draw[fill=red]  (0,0) circle (1.5pt) ;
\node[rectangle,draw, right, baseline] (x) at (0,1) {\LARGE\color{black!30}{before}\relax};
\draw[fill=red]  (0,1) circle (1.5pt) ;
\node[rectangle,draw, right=1sp] at (0,0) {\LARGE\color{black!20} \so\chaptername\relax};

\node[rectangle,draw, color=white, below right, fill=blue!50, text=white] at +(\textwidth+3cm,0) {\scalebox{2}{\HHUGE \thechapter}};
\draw[fill=red]  (0,-3) circle (1.5pt) ;
\node[rectangle, draw, text width=9cm,below right, yshift=-1pt] at (0,-3) {%
         \HUGE Title format\kern0ptt\kern0pting \Large \lorem }; 
\node at (12.5,-9) {\includegraphics[width=7cm]{./images/fashion.jpg}};
\hook
\end{tikzpicture}

\tikzrule 
\egroup

\lipsum
\bigskip

\thispagestyle{empty}

\begin{tikzpicture}
\draw[fill=black!20] (0,0) rectangle (1+1,2) rectangle ++(1+1,-2) {};


\end{tikzpicture}

\end{document}

http://tex.stackexchange.com/questions/48473/best-way-to-give-sources-of-images-used-in-a-beamer-presentation/48479#48479